%!TeX root=../pridetop.tex
\chapter[Chapter \thechapter]{}
	
\begin{figure}[t!]
	\centering
\includegraphics[width=.8\linewidth]{17delighted}
\captionlistentry{Delighted to see their dear friend again}
\end{figure}	
	
\lettrine[lines=6,image=true]{initials/chap17e}{lizabeth}  related to Jane, the next day, what had passed between Mr Wickham and herself. Jane listened with astonishment and concern: she knew not how to believe that Mr Darcy could be so unworthy of Mr Bingley's regard; and yet it was not in her nature to question the veracity of a young man of such amiable appearance as Wickham. The possibility of his having really endured such unkindness was enough to interest all her tender feelings; and nothing therefore remained to be done but to think well of them both, to defend the conduct of each, and throw into the account of accident or mistake whatever could not be otherwise explained.

<They have both,> said she, <been deceived, I dare say, in some way or other, of which we can form no idea. Interested people have perhaps misrepresented each to the other. It is, in short, impossible for us to conjecture the causes or circumstances which may have alienated them, without actual blame on either side.>

<Very true, indeed; and now, my dear Jane, what have you got to say in behalf of the interested people who have probably been concerned in the business? Do clear \textit{them}, too, or we shall be obliged to think ill of somebody.>

<Laugh as much as you choose, but you will not laugh me out of my opinion. My dearest Lizzy, do but consider in what a disgraceful light it places Mr Darcy, to be treating his father's favourite in such a manner,—one whom his father had promised to provide for. It is impossible. No man of common humanity, no man who had any value for his character, could be capable of it. Can his most intimate friends be so excessively deceived in him? Oh no.>

<I can much more easily believe Mr Bingley's being imposed on than that Mr Wickham should invent such a history of himself as he gave me last night; names, facts, everything mentioned without ceremony. If it be not so, let Mr Darcy contradict it. Besides, there was truth in his looks.>

<It is difficult, indeed—it is distressing. One does not know what to think.>

<I beg your pardon;—one knows exactly what to think.>

But Jane could think with certainty on only one point,—that Mr Bingley, if he \textit{had} been imposed on, would have much to suffer when the affair became public.

The two young ladies were summoned from the shrubbery, where this conversation passed, by the arrival of some of the very persons of whom they had been speaking; Mr Bingley and his sisters came to give their personal invitation for the long expected ball at Netherfield, which was fixed for the following Tuesday. The two ladies were delighted to see their dear friend again, called it an age since they had met, and repeatedly asked what she had been doing with herself since their separation. To the rest of the family they paid little attention; avoiding Mrs Bennet as much as possible, saying not much to Elizabeth, and nothing at all to the others. They were soon gone again, rising from their seats with an activity which took their brother by surprise, and hurrying off as if eager to escape from Mrs Bennet's civilities.



The prospect of the Netherfield ball was extremely agreeable to every female of the family. Mrs Bennet chose to consider it as given in compliment to her eldest daughter, and was particularly flattered by receiving the invitation from Mr Bingley himself, instead of a ceremonious card. Jane pictured to herself a happy evening in the society of her two friends, and the attentions of their brother; and Elizabeth thought with pleasure of dancing a great deal with Mr Wickham, and of seeing a confirmation of everything in Mr Darcy's look and behaviour. The happiness anticipated by Catherine and Lydia depended less on any single event, or any particular person; for though they each, like Elizabeth, meant to dance half the evening with Mr Wickham, he was by no means the only partner who could satisfy them, and a ball was, at any rate, a ball. And even Mary could assure her family that she had no disinclination for it.

<While I can have my mornings to myself,> said she, <it is enough. I think it is no sacrifice to join occasionally in evening engagements. Society has claims on us all; and I profess myself one of those who consider intervals of recreation and amusement as desirable for everybody.>

Elizabeth's spirits were so high on the occasion, that though she did not often speak unnecessarily to Mr Collins, she could not help asking him whether he intended to accept Mr Bingley's invitation, and if he did, whether he would think it proper to join in the evening's amusement; and she was rather surprised to find that he entertained no scruple whatever on that head, and was very far from dreading a rebuke, either from the Archbishop or Lady Catherine de Bourgh, by venturing to dance.

<I am by no means of opinion, I assure you,> said he, <that a ball of this kind, given by a young man of character, to respectable people, can have any evil tendency; and I am so far from objecting to dancing myself, that I shall hope to be honoured with the hands of all my fair cousins in the course of the evening; and I take this opportunity of soliciting yours, Miss Elizabeth, for the two first dances especially; a preference which I trust my cousin Jane will attribute to the right cause, and not to any disrespect for her.>

Elizabeth felt herself completely taken in. She had fully proposed being engaged by Wickham for those very dances; and to have Mr Collins instead!—her liveliness had been never worse timed. There was no help for it, however. Mr Wickham's happiness and her own was perforce delayed a little longer, and Mr Collins's proposal accepted with as good a grace as she could. She was not the better pleased with his gallantry, from the idea it suggested of something more. It now first struck her, that \textit{she} was selected from among her sisters as worthy of being the mistress of Hunsford Parsonage, and of assisting to form a quadrille table at Rosings, in the absence of more eligible visitors. The idea soon reached to conviction, as she observed his increasing civilities towards herself, and heard his frequent attempt at a compliment on her wit and vivacity; and though more astonished than gratified herself by this effect of her charms, it was not long before her mother gave her to understand that the probability of their marriage was exceedingly agreeable to \textit{her}. Elizabeth, however, did not choose to take the hint, being well aware that a serious dispute must be the consequence of any reply. Mr Collins might never make the offer, and, till he did, it was useless to quarrel about him.

If there had not been a Netherfield ball to prepare for and talk of, the younger Miss Bennets would have been in a pitiable state at this time; for, from the day of the invitation to the day of the ball, there was such a succession of rain as prevented their walking to Meryton once. No aunt, no officers, no news could be sought after; the very shoe-roses for Netherfield were got by proxy. Even Elizabeth might have found some trial of her patience in weather which totally suspended the improvement of her acquaintance with Mr Wickham; and nothing less than a dance on Tuesday could have made such a Friday, Saturday, Sunday, and Monday endurable to Kitty and Lydia.